%!TEX root = ../dissertation.tex
\chapter{Introduction}\label{chapt:intro}

% =====================================================================
%  DRAFT v1 (Asma) -- full prose. Search "HIGHLIGHT" for things only you
%  can decide/insert (figures, exact citations, personal framing).
%  Target length ~8 pages.
% =====================================================================

\section{Motivation}\label{sect:intro_motivation}

Searching a large image collection for one specific picture is a task that
text alone, or an example image alone, often cannot express. A shopper who has
seen \emph{this exact sneaker, but in the high-top version}, a collector looking
for \emph{this porcelain vase, but with the floral pattern in blue}, or a tourist
who recalls \emph{this landmark, rendered as a vintage engraving}, all have a
target in mind that is naturally described as an existing image together with a
desired change. Neither modality suffices on its own: the words do not pin down
which sneaker or which vase, and the example image does not capture the change the
user wants. This
is the setting of \emph{composed image retrieval} (CIR): the query is a pair
$(\qv, \qt)$ consisting of a reference image $\qv$ and a short text $\qt$ that
specifies how the desired result differs from it, and the system must return,
from a database $\DB$, the image the user actually has in mind.

Composed image retrieval has attracted growing attention because it combines
the complementary strengths of the two most common query modalities. The
reference image grounds the search in concrete visual appearance, while the
text supplies the relational or attribute-level change that an image alone
cannot convey~\cite{vo2019tirg}. Applications range from e-commerce product
search and fashion recommendation~\cite{wu2021fashioniq} to general visual
search on real-world photographs~\cite{liu2021cirr}. The central difficulty is
that the modification text does not describe the reference image; it describes
the \emph{difference} between the reference and an image that is not available
at query time, which the model must imagine and then match against the
database.

\section{Instance-Level Composed Image Retrieval}\label{sect:intro_icir}

Most CIR benchmarks are \emph{category-level}: a retrieved image is considered
correct if it belongs to the same broad category as the intended target (for
example, ``a red dress'' rather than \emph{the} specific dress the user
wanted). Recent work argues that this protocol overstates real performance,
because it rewards coarse semantic matches and tolerates retrieving a
different object of the same kind. The \emph{instance-level} composed image
retrieval benchmark, i-CIR~\cite{psomas2025icir}, was introduced to close this
gap. In i-CIR the ground truth is defined at the level of a specific object
\emph{instance}: a retrieved image is correct only if it depicts the same
instance, modified as the query text requests.

% HIGHLIGHT: confirm the exact dataset statistics below against the i-CIR
% paper / your data card before final submission.
This makes the task substantially harder in two ways. First, the database is
large and open-domain (752{,}092 database images and 1{,}883 queries in the
full benchmark), so the correct target competes with a very large pool of
distractors. Second, and more importantly, i-CIR deliberately populates each
query's pool with \emph{hard negatives}: images of the same instance under a
\emph{different} modification, visually similar images that do \emph{not}
satisfy the text, and images that satisfy the text but show a \emph{different}
instance. A method that latches onto only the reference image, or only the
text, is therefore systematically punished. Success requires jointly
respecting both the visual identity of the instance and the precise change
demanded by the text.

\section{Zero-Shot Retrieval and the Role of Language}\label{sect:intro_gap}

A natural way to attack CIR is to train a fusion network that maps the pair
$(\qv, \qt)$ to a single query embedding using triplet supervision on a CIR
training set~\cite{vo2019tirg}. The obstacle is the supervision itself. Training
such a model requires a large collection of triplets of the form (reference
image, modification text, positive image), and these are difficult to obtain at
scale. In practice they come from one of three sources, each with a drawback.
They may be \emph{manually} annotated, which limits datasets to a few tens of
thousands of carefully labelled triplets~\cite{vo2019tirg,wu2021fashioniq,liu2021cirr,baldrati2023circo};
they may be \emph{crawled} from the web by pairing images that co-occur on the
same page~\cite{zhang2024magiclens,ventura2024covr}; or they may be
\emph{synthesised} by a generative model~\cite{gu2024compodiff}. The automatic
routes scale, but at a cost in quality, and in every case the diversity of visual
object types and textual modifications remains far narrower than the variety a
vision--language model already sees during pre-training, which limits how well a
trained fusion model generalises beyond its training distribution. Instance-level
annotation is more expensive still, and the i-CIR benchmark is explicitly
designed for evaluation rather than for large-scale training.

These considerations motivate the \emph{zero-shot} setting~\cite{saito2023pic2word,baldrati2023circo},
in which no CIR-specific training is performed and retrieval is built entirely on
\emph{off-the-shelf}, frozen vision--language models such as
CLIP~\cite{radford2021clip} and SigLIP~\cite{zhai2023siglip,tschannen2025siglip2}.
A zero-shot method inherits the broad coverage of the pre-trained encoders
directly, needs no triplet collection, and is not tied to the distribution of any
one CIR training set -- which is precisely the setting this thesis adopts.

The i-CIR paper accompanies the benchmark with a strong zero-shot baseline,
\textsc{Basic}, which combines a frozen vision--language backbone with a
sequence of feature post-processing steps (feature centering, dimensionality
reduction, query expansion, and a contextualization of the modification text)
to fuse the image and text similarities~\cite{psomas2025icir}. \textsc{Basic}
shows that careful use of frozen features alone can be ``surprisingly strong''.
Yet it also exposes the core weakness of treating the modification text in
isolation: a short phrase such as \emph{``make it red''} is a poor standalone
text query, because the encoder has no access to \emph{what} should become red.
The text only makes sense \emph{relative} to the reference image.

This thesis builds on the observation that language can serve as an explicit
bridge between the two modalities. Instead of feeding the encoder the bare
modification text, we use a multimodal large language model (MLLM) to read the
reference image and the modification \emph{together} and to write a short
description of the \emph{imagined target} -- effectively performing the edit in
language space. This \emph{target caption} is a self-contained text query that
the frozen text encoder can embed well, and it captures the composition that
the raw modification text cannot.

\section{Contributions}\label{sect:intro_contrib}

The contributions of this thesis are the following:

\begin{itemize}
  \item \textbf{An MLLM-generated target caption as a third retrieval modality.}
        We propose to generate, for each query, a caption of the intended
        target image with a frozen MLLM (InternVL-3.5-8B~\cite{wang2025internvl35}), and to use the
        embedding of this caption as an additional, text-side retrieval signal
        alongside the reference-image and modification-text signals
        (Chapter~\ref{chapt:method}).
  \item \textbf{A triplet late-fusion scoring function.} We rank database
        images by the clamped product of three similarities -- image, text,
        and caption -- which acts as a soft logical \textsc{and}: a candidate
        must agree with all three signals to score highly.
  \item \textbf{A backbone comparison on the full i-CIR benchmark.} We embed
        the full database (752K images) with four off-the-shelf backbones
        (CLIP ViT-L/14~\cite{radford2021clip}, CLIP ViT-H/14~\cite{ilharco2021openclip},
        SigLIP ViT-L/16~\cite{zhai2023siglip}, and SigLIP2 g/16-384~\cite{tschannen2025siglip2})
        and report apples-to-apples results, identifying SigLIP2 as the
        strongest backbone for our pipeline.
  \item \textbf{Ablations isolating each component.} We quantify the separate
        contributions of the caption, of feature centering, and of the
        remaining \textsc{Basic} preprocessing, and show that the caption is
        the single largest contributor -- adding it to a plain text$\times$image
        baseline reaches, and with a strong backbone exceeds, the full
        \textsc{Basic} preprocessing stack (Chapter~\ref{chapt:results}).
\end{itemize}

% HIGHLIGHT (optional but recommended): add a small "teaser" figure here --
% one query (reference image + text), the generated caption, and the top-1
% retrieved image -- so the reader sees the idea on page 1.
% \begin{figure}[t]
%   \centering
%   \includegraphics[width=\textwidth]{figures/teaser.png}
%   \caption{...}\label{fig:teaser}
% \end{figure}

\section{Outline}\label{sect:intro_outline}

This thesis is organised as follows.
Chapter~\ref{chapt:background} provides the theoretical background on image
retrieval, contrastive vision--language pretraining, and multimodal large
language models.
Chapter~\ref{chapt:related} reviews related work on composed image retrieval and
zero-shot CIR.
Chapter~\ref{chapt:benchmark} introduces the i-CIR benchmark used throughout: the
problem formulation, the dataset and its statistics, the evaluation protocol, and
the \textsc{Basic} baseline that this thesis builds on.
Chapter~\ref{chapt:method} describes the proposed methodology: the retrieval
pipeline, the vision--language backbones, the target-caption generation step, and
the triplet fusion.
Chapter~\ref{chapt:results} reports and discusses the experimental results on
the full i-CIR dataset and the ablation studies on the development subset.
Chapter~\ref{chapt:conclusion} draws the conclusions and outlines future work.
